% Compile with xelatex
\documentclass[11pt,a4paper]{article}

\usepackage{fontspec}
\setmainfont{Malgun Gothic}[Ligatures=TeX]
\setsansfont{Malgun Gothic}[Ligatures=TeX]

\usepackage[top=2.2cm, bottom=2.2cm, left=2.5cm, right=2.5cm]{geometry}
\usepackage{booktabs,longtable,array,multirow}
\usepackage{graphicx}
\usepackage{float}
\usepackage{amsmath}
\usepackage[hidelinks]{hyperref}

\setlength{\parindent}{0pt}
\setlength{\parskip}{6pt}

\newcolumntype{R}{>{\raggedleft\arraybackslash}p{1.4cm}}

% ─────────────────────────────────────────────────────────────
\begin{document}

\begin{center}
{\Large\textbf{Preliminary Analysis Report}}\\[6pt]
{\large Cross-Cultural Perception of Korean Tourist Attractions}\\[4pt]
{\normalsize 8 places (Seoul + Busan) $\cdot$ 496 reviews $\cdot$ October 2026}
\end{center}

\vspace{8pt}
\noindent\rule{\linewidth}{0.4pt}

% ── 1. CORPUS ────────────────────────────────────────────────
\section*{1. Corpus Summary}

A bilingual corpus of 496 reviews was constructed from 8 tourist destinations
in Seoul and Busan. Each place is represented by approximately 30 Korean-language
(domestic) and 30 English-language (foreign) reviews, collected from
culturally segmented online review platforms.

\begin{table}[H]
\centering
\small
\begin{tabular}{llcc}
\toprule
\textbf{City} & \textbf{Place} & \textbf{KO reviews} & \textbf{EN reviews} \\
\midrule
Seoul  & Gyeongbokgung (경복궁)         & 30 & 30 \\
Seoul  & Changdeokgung (창덕궁)         & 30 & 30 \\
Seoul  & Bukchon Village (북촌한옥마을)   & 32 & 30 \\
Seoul  & Gwangjang Market (광장시장)     & 30 & 30 \\
Seoul  & Hongdae Street (홍대거리)       & 34 & 30 \\
Seoul  & Lotte World (롯데월드)          & 34 & 30 \\
Busan  & Gwangalli Beach (광안리해수욕장) & 34 & 30 \\
Busan  & Haeundae Beach (해운대해수욕장)  & 32 & 30 \\
\midrule
       & \textbf{Total}                 & \textbf{256} & \textbf{240} \\
\bottomrule
\end{tabular}
\caption{Corpus composition by place and language group.}
\end{table}

% ── 2. SENTIMENT ─────────────────────────────────────────────
\section*{2. Sentiment Analysis}

Pre-trained sentiment classifiers were applied to each language group separately.
The following table shows the mean sentiment score (range 0--1) and the
result of a one-sided Mann-Whitney U test (H$_1$: KO $>$ EN).

\begin{table}[H]
\centering
\small
\begin{tabular}{lRRRl}
\toprule
\textbf{Place} & \textbf{KO} & \textbf{EN} & \textbf{p-value} & \textbf{Sig.} \\
\midrule
Gyeongbokgung    & 0.597 & 0.673 & 0.7855 & n.s. \\
Changdeokgung    & 0.524 & 0.668 & 0.9253 & n.s. \\
Bukchon Village  & 0.408 & 0.469 & 0.5782 & n.s. \\
Gwangjang Market & 0.282 & 0.455 & 0.6078 & n.s. \\
Hongdae Street   & 0.621 & 0.419 & 0.0509 & n.s. \\
Lotte World      & 0.826 & 0.597 & 0.0362 & * \\
Gwangalli Beach  & 0.745 & 0.775 & 0.9648 & n.s. \\
Haeundae Beach   & 0.643 & 0.758 & 0.9314 & n.s. \\
\midrule
\textbf{Overall} & \textbf{0.587} & \textbf{0.602} & \textbf{0.6559} & n.s. \\
\bottomrule
\end{tabular}
\caption{Mean sentiment and Mann-Whitney U test results. * $p < 0.05$.}
\end{table}

Lotte World is the only place with a statistically significant gap
(KO $>$ EN, $p = 0.036$). Hongdae is borderline ($p = 0.051$).
Gwangjang Market has the lowest domestic sentiment (0.282),
suggesting more critical local reviewers for food-related destinations.

% ── 3. CROSS-LINGUAL SIMILARITY ─────────────────────────────
\section*{3. Cross-lingual Semantic Similarity}

Multilingual sentence embeddings (multilingual-e5-base) project both
Korean and English reviews into a shared semantic space. Per-site cosine
similarity between group centroids quantifies experiential convergence.

\begin{table}[H]
\centering
\small
\begin{tabular}{lRRR}
\toprule
\textbf{Place} & \textbf{KO-EN sim} & \textbf{KO intra} & \textbf{EN intra} \\
\midrule
Gyeongbokgung    & 0.910 & 0.879 & 0.860 \\
Changdeokgung    & 0.916 & 0.872 & 0.859 \\
Bukchon Village  & 0.895 & 0.873 & 0.878 \\
Gwangjang Market & 0.931 & 0.869 & 0.857 \\
Hongdae Street   & 0.922 & 0.870 & 0.870 \\
Lotte World      & 0.929 & 0.866 & 0.852 \\
Gwangalli Beach  & 0.909 & 0.892 & 0.889 \\
Haeundae Beach   & 0.896 & 0.895 & 0.892 \\
\bottomrule
\end{tabular}
\caption{Cross-lingual cosine similarity and intra-group coherence.}
\end{table}

Cross-group similarity is consistently high (0.89--0.93) and exceeds
intra-group similarity in all cases, indicating that both groups discuss
similar aspects at each location. Gwangjang Market shows the highest
alignment (0.931), while Haeundae Beach shows the most divergence (0.896).

% ── 4. MOTIVATION ────────────────────────────────────────────
\section*{4. Travel Motivation (Zero-shot Classification)}

A zero-shot classification using cosine similarity in the embedding
space assigned each review to one of five motivation categories.

\begin{table}[H]
\centering
\small
\begin{tabular}{lRR}
\toprule
\textbf{Motivation} & \textbf{KO (\%)} & \textbf{EN (\%)} \\
\midrule
Photography \& Visual Experience & 37.5 & 37.1 \\
Nature \& Scenery                & 25.0 & 21.7 \\
Cultural Heritage \& History     & 16.2 & 15.8 \\
Food \& Gastronomy               & 16.2 & 15.8 \\
Shopping \& Entertainment        &  5.0 &  9.6 \\
\bottomrule
\end{tabular}
\caption{Overall motivation distribution by language group.}
\end{table}

Both groups are dominated by Photography/Visual ($\sim$37\%) and
Nature/Scenery ($\sim$21--25\%). The only notable gap is in
Shopping \& Entertainment, where English reviewers show nearly
twice the proportion (9.6\% vs 5.0\%).

% ── 5. LEXICAL DIVERSITY ────────────────────────────────────
\section*{5. Lexical Diversity (TTR)}

Type-Token Ratio (TTR) measures lexical diversity per review.
English reviewers show consistently higher TTR across all places
(mean EN = 0.937, mean KO = 0.873), suggesting more varied vocabulary
usage. This may reflect differences in review writing conventions
between platforms or the linguistic properties of each language.

% ── 6. TF-IDF KEYWORDS ──────────────────────────────────────
\section*{6. TF-IDF Keyword Contrast}

Top distinctive keywords per group (by mean TF-IDF score):

\begin{table}[H]
\centering
\small
\begin{tabular}{clcl}
\toprule
\textbf{Rank} & \textbf{Korean} & & \textbf{English} \\
\midrule
1 & 시간 (time)     & & beach \\
2 & 경복궁 (Gyeongb.)& & palace \\
3 & 해운대 (Haeundae)& & worth \\
4 & 바다 (sea)      & & area \\
5 & 홍대 (Hongdae)  & & experience \\
6 & 한옥 (hanok)    & & food \\
7 & 사진 (photo)    & & street \\
8 & 풍경 (scenery)  & & traditional \\
\bottomrule
\end{tabular}
\caption{Top TF-IDF keywords per language group.}
\end{table}

Korean reviewers emphasize temporal and logistical aspects (시간, time),
while English reviewers focus on evaluative language (worth, experience).

% ── 7. LDA TOPICS ────────────────────────────────────────────
\section*{7. LDA Topic Modeling}

Five topics were extracted for each language group via Latent Dirichlet
Allocation.

\textbf{English Topics:}
\begin{enumerate}\itemsep2pt
  \item beach, clean, enjoy, rides, people, bridge
  \item palace, garden, secret, changdeokgung, traditional
  \item area, fun, beach, people, food, street, local
  \item beach, busan, walking, haeundae, atmosphere, ocean
  \item beach, food, street, night, area, busan, walk
\end{enumerate}

\textbf{Korean Topics:}
\begin{enumerate}\itemsep2pt
  \item 창덕궁, 시장, 광장시장, 전통, 다양, 빈대떡
  \item 홍대, 거리, 시간, 다양, 북촌, 양평, 골목
  \item 경복궁, 해운대, 사진, 드론, 광안리, 외국인
  \item 창덕궁, 궁궐, 조선, 자연, 경복궁, 후원
  \item 한옥, 바다, 풍경, 부산, 해변, 마을, 대교
\end{enumerate}

% ── 8. OBSERVATIONS ──────────────────────────────────────────
\section*{8. Preliminary Observations}

\begin{itemize}\itemsep3pt
  \item No overall significant sentiment difference between domestic and
        foreign reviewers ($p = 0.656$); differences are place-specific.
  \item Lotte World shows the strongest cultural gap: Korean reviewers
        are significantly more positive ($p = 0.036$).
  \item High cross-lingual similarity (0.89--0.93) suggests both groups
        discuss similar aspects, but with different emphasis and vocabulary.
  \item Photography/Visual dominates both groups' motivations ($\sim$37\%),
        while Shopping/Entertainment is more salient for foreign visitors.
  \item Data collection for 2 additional places (인사동길, N서울타워) is pending.
\end{itemize}

\vspace{8pt}
\noindent\rule{\linewidth}{0.4pt}
{\small Hoseo University \hfill October 2026}

\end{document}
